\chapter{Complexity — Poly Time}
%%% generated by the blueprint pipeline from TCSlib.Complexity.ClassNP.PolyTime
%%%%%%%%%%%%%%%%%%%%%%%%%%%%%%%%%%%%%%%%%%%%%%%%%%%%%%%%%%%%%%%%%%%%%%%%%%%%%%%
\section{Overview}

This module introduces the class $\mathrm{FP}$ of polynomial-time computable functions on
bit strings that underlies every Karp reduction of [AB09, ch.~2]: a function
$f : \{0,1\}^* \to \{0,1\}^*$ is polynomial-time computable when some finite binary Turing
machine computes it within $C(n+1)^c$ steps on inputs of length $n$. It fixes this
polynomial normal form and proves the basic closure calculus used by the chapter's
reductions: the identity is polynomial-time computable, polynomial-time functions have
polynomially bounded output length, and the class is closed under composition
[AB09, Theorem~2.8].

%%%%%%%%%%%%%%%%%%%%%%%%%%%%%%%%%%%%%%%%%%%%%%%%%%%%%%%%%%%%%%%%%%%%%%%%%%%%%%%
\section{Declarations}

\begin{definition}[Polynomially bounded function]
\lean{Complexity.PolyBound}
\label{Complexity.PolyBound}
\leanok
A function $p : \mathbb{N} \to \mathbb{N}$ is \emph{polynomially bounded} if there are
constants $C, c \in \mathbb{N}$ such that $p(n) \le C(n+1)^c$ for every $n \in \mathbb{N}$.
The function $p$ itself need be neither monotone nor computable; only the majorant
$C(n+1)^c$ is.
\end{definition}

\begin{definition}[Polynomial-time computable function]
\lean{Complexity.PolyTimeComputable}
\label{Complexity.PolyTimeComputable}
\leanok
\uses{Turing.FinTM, Turing.FinTM.ComputesFunInTime}
A function $f : \{0,1\}^* \to \{0,1\}^*$ on bit strings is \emph{polynomial-time computable}
if there exist a Turing machine $M$ over the binary alphabet with finitely many states and
constants $C, c \in \mathbb{N}$ such that, on every input $x$ of length $n$, $M$ halts within
$C(n+1)^c$ steps with $f(x)$ written on its output tape. These functions form the class
$\mathrm{FP}$ implicit throughout [AB09, ch.~2].
\end{definition}

\begin{theorem}[The identity is polynomial-time computable]
\lean{Complexity.polyTimeComputable_id}
\label{Complexity.polyTimeComputable_id}
\leanok
\uses{Complexity.PolyTimeComputable, Turing.FinTM.ComputesInTime.mono, Turing.FinTM.computesFunInTime_id}
\difficulty{3}
The identity function $x \mapsto x$ on bit strings is polynomial-time computable: some finite
binary Turing machine outputs $x$ within $C(n+1)^c$ steps on every input $x$ of length $n$, for
some constants $C, c \in \mathbb{N}$.
\end{theorem}

\begin{theorem}[Polynomial-time functions have polynomial output length]
\lean{Complexity.PolyTimeComputable.output_length_le}
\label{Complexity.PolyTimeComputable.output_length_le}
\leanok
\uses{Complexity.PolyTimeComputable, Turing.FinTM.computesInTime_iff, Turing.MultiTapeTM.output_length_le}
\difficulty{3}
Let $f : \{0,1\}^* \to \{0,1\}^*$ be polynomial-time computable, i.e.\ some finite binary
Turing machine computes $f$ within $C_0(n+1)^{c_0}$ steps on every input of length $n$, for
some constants $C_0, c_0 \in \mathbb{N}$. Then there are constants $C, c \in \mathbb{N}$ such
that $|f(x)| \le C(|x|+1)^c$ for every bit string $x$.
\end{theorem}

\begin{lemma}[Polynomial bound for a composed time budget]
\lean{Complexity.comp_time_bound}
\label{Complexity.comp_time_bound}
\leanok
\difficulty{4}
For all natural numbers $a, C, c, C', c', n$,
\[
  a\bigl(C(n+1)^c + C'\,(C(n+1)^c + 1)^{c'} + 1\bigr)
  \;\le\;
  a\bigl(C + C'(C+1)^{c'} + 1\bigr)\,(n+1)^{\max(c,\, c c')}.
\]
\end{lemma}

\begin{theorem}[Polynomial-time functions are closed under composition]
\lean{Complexity.PolyTimeComputable.comp}
\label{Complexity.PolyTimeComputable.comp}
\leanok
\uses{Complexity.PolyTimeComputable, Complexity.comp_time_bound, Turing.FinTM.ComputesInTime.mono, Turing.FinTM.computesFunInTime_comp}
\difficulty{3}
Let $f, g : \{0,1\}^* \to \{0,1\}^*$ be polynomial-time computable, where a function is
polynomial-time computable if some finite binary Turing machine computes it within $C(n+1)^c$
steps on every input of length $n$, for some constants $C, c \in \mathbb{N}$. Then the
composition $g \circ f$, mapping $x$ to $g(f(x))$, is also polynomial-time computable
[AB09, proof of Theorem~2.8].
\end{theorem}
